\section{Introduction}\label{sec:intro}

Vehicle trajectories are collected at scale by roadside sensors, connected-vehicle platforms and aftermarket loggers, usually without a reliable record of the vehicle that produced them. Many uses need it: traffic analytics by vehicle class, calibration of microscopic simulation models whose car-following parameters differ between a compact car and a three-row SUV, and infrastructure studies of how often heavy or electric vehicles use a facility. Registration records, licence plates or camera images can fill the gap but raise privacy and availability problems; a trajectory that reveals its vehicle type by itself would be more privacy-friendly.

Motion should carry this information: wheelbase, steering ratio and mass distribution set the yaw and curvature response to steering, and the powertrain shapes launches and lift-off deceleration, so a battery-electric vehicle (BEV) braking largely through its motor should leave a different speed trace from an internal-combustion-engine (ICE) vehicle. We pose the question as it arises in deployment, for car models never seen in training, since otherwise every new model would need retraining. From 60-s windows of trajectory and controller area network (CAN) signals (among them speed, acceleration, steering, yaw rate, curvature, lateral acceleration, pedal and automation flags, and lead-vehicle distance, relative speed, headway and time-to-collision; Section~\ref{sec:dataset} lists all channels) we predict a size class (SMALL, MID or LARGE, cut from the vehicle footprint) and the powertrain (BEV vs ICE).

Three things make this hard. First, the size label is a hard cut through a continuous quantity (footprints of 4.43 and 4.88 m$^2$ separate the classes), so a car model near a cut is hard to place even when its physics is read well. Second, a driving automation system (DAS) is engaged in 47.9\%
of the samples, so part of the measured behaviour is a controller signature that may not transfer between car models. Third, the fleet is uneven (174 car models from 38 brands with widely differing hours), so a held-out model in a footprint gap of the training fleet has no neighbour to be read against.

We therefore release the data with fixed labels and splits, benchmark sequence models and feature pipelines under one protocol, run a pre-registered method study on model-disjoint folds, and ask, without training anything, where the size ceiling lies and why. The contributions are:

\begin{itemize}
\item A public, privacy-preserving trajectory and CAN dataset (9,062 drives, about 3,950 h, 1,431 pseudonymous vehicles) with fixed footprint-size and powertrain labels, held-out-car-model (headline) and held-out-device splits, and window-, drive- and vehicle-level protocols with cluster-bootstrap intervals over car models (Section~\ref{sec:dataset}).
\item A benchmark of four sequence architectures, hand-crafted statistics and physical features proposed by a large language model (LLM) under those protocols, showing that architecture is not the bottleneck (all deep models within about 2 points) and drive-level aggregation adds 7--12 points (Table~\ref{tab:main}).
\item A pre-registered method study on model-disjoint internal folds (Sections~\ref{sec:methods} and \ref{sec:experiments}): car-model-balanced sampling, auxiliary physical targets, physics-safe augmentation, per-car-model loss re-weighting (size) and DAS down-weighting (BEV vs ICE) raise the drive-level result by 4--6 points (size) and 2--4 points (BEV vs ICE); an episode-centric multiple-instance model over driving transients matches but does not beat the window recipes on the test models.
\item Zero-training ceiling analyses (per-car-model majority, footprint-noise and curb-mass oracles, evidence sufficiency, separability of held-out models) that quantify why size cannot reach 90 on unseen car models under a footprint cut, and what would make it reachable (Figs.~\ref{fig:evidence} and \ref{fig:ceiling}).
\end{itemize}
